% Define important features for the document
\documentclass[12pt,oneside]{article}
\setlength{\parindent}{0pt}
\setlength{\parskip}{4pt}

% Import packages
\usepackage[margin = 1in]{geometry}
\usepackage{amsfonts, amsmath, amssymb, amsthm, calc, centernot, color, csquotes, dsfont, enumitem, environ, fancyhdr, fontawesome5, graphicx, hyperref, listings, marvosym, mathrsfs, mathtools, tcolorbox, titlesec, upquote, xcolor}
\usepackage[flushmargin,hang]{footmisc}
\tcbuselibrary{breakable}

% Define toggles
\newif\ifsolutions
\solutionstrue % Toggle for instructor version
% \solutionsfalse % Toggle for student version

% Define size and spacing of section and subsection
\titleformat{\section}
  {\normalfont\large\bfseries}
  {\thesection}{1em}{}
\titleformat{\subsection}
  {\normalfont\normalsize\bfseries}
  {\thesubsection}{1em}{}
\titlespacing*{\section}{0pt}{6pt}{6pt} % ({left}{before}{after})
\titlespacing*{\subsection}{0pt}{6pt}{6pt} % ({left}{before}{after})

% Define commands
\newcommand{\indep}{\perp\!\!\!\perp}
\newcommand{\notindep}{\mathrel{\not\!\perp\!\!\!\perp}}
\newcommand{\E}{\mathbb{E}}
\newcommand{\Var}{\text{Var}}
\newcommand{\Cov}{\text{Cov}}
\newcommand{\Corr}{\text{Corr}}
\newcommand{\SD}{\text{SD}}
\newcommand{\SE}{\text{SE}}
\newcommand{\Bias}{\text{Bias}}
\newcommand{\MSE}{\text{MSE}}
\newcommand{\Risk}{\text{Risk}}
\newcommand{\Like}{\mathcal{L}}
\newcommand{\boot}{\text{boot}}
\newcommand{\strat}{\text{strat}}
\newcommand{\supp}{\text{supp}}
\newcommand{\MLE}{\text{MLE}}
\newcommand{\MoM}{\text{MoM}}
\newcommand{\MOM}{\text{MOM}}
\newcommand{\OLS}{\text{OLS}}
\newcommand{\WLS}{\text{WLS}}
\newcommand{\LS}{\text{LS}}
\newcommand{\MAP}{\text{MAP}}
\newcommand{\MAE}{\text{MAE}}
\newcommand{\LP}{\text{LP}}
\newcommand{\HT}{\text{HT}}
\newcommand{\JS}{\text{JS}}
\newcommand{\RB}{\text{RB}}
\newcommand{\UB}{\text{UB}}
\newcommand{\DIM}{\text{DIM}}
\newcommand{\PATE}{\text{PATE}}
\newcommand{\SATE}{\text{SATE}}
\newcommand{\VIF}{\text{VIF}}
\newcommand{\GVIF}{\text{GVIF}}
\newcommand{\AIC}{\text{AIC}}
\newcommand{\BIC}{\text{BIC}}
\newcommand{\LASSO}{\text{LASSO}}
\newcommand{\PSS}{\text{PSS}}
\newcommand{\CV}{\text{CV}}
\newcommand{\logit}{\text{logit}}
\newcommand{\odds}{\text{odds}}
\newcommand{\OR}{\text{OR}}
\newcommand{\N}{\mathcal{N}}
\newcommand{\Bern}{\text{Bern}}
\newcommand{\Bin}{\text{Bin}}
\newcommand{\Beta}{\text{Beta}}
\newcommand{\Gam}{\text{Gamma}}
\newcommand{\Expo}{\text{Expo}}
\newcommand{\Pois}{\text{Pois}}
\newcommand{\Unif}{\text{Unif}}
\newcommand{\DUnif}{\text{DUnif}}
\newcommand{\Geom}{\text{Geom}}
\newcommand{\FS}{\text{FS}}
\newcommand{\NBin}{\text{NBin}}
\newcommand{\HGeom}{\text{HGeom}}
\newcommand{\Mult}{\text{Mult}}
\newcommand{\MVN}{\text{MVN}}
\newcommand{\IG}{\text{IG}}
\newcommand{\Tweedie}{\text{Tweedie}}
\newcommand{\simiid}{\overset{\text{i.i.d.}}{\sim}}
\newcommand{\simind}{\overset{\text{ind.}}{\sim}}
\newcommand{\simnull}{\overset{H_0}{\sim}}
\newcommand{\simapprox}{\overset{\cdot}{\sim}}
\newcommand{\simnullapprox}{\overset{H_0}{\simapprox}}
\newcommand{\R}{\mathbb{R}}
\newcommand{\Z}{\mathbb{Z}}
\newcommand{\Q}{\mathbb{Q}}
\newcommand{\C}{\mathbb{C}}
\newcommand{\I}{\mathbb{I}}
\newcommand{\indic}{\mathds{1}}
\newcommand{\iid}{\text{i.i.d.}}
\newcommand{\Ind}{\text{Ind}}
\newcommand{\SSq}{\text{SSq}}
\newcommand{\trace}{\text{trace}}
\newcommand{\rank}{\text{rank}}
\newcommand{\diag}{\text{diag}}
\newcommand{\adj}{\text{adj}}
\newcommand{\SSE}{\text{SSE}}
\newcommand{\SST}{\text{SST}}
\newcommand{\SSR}{\text{SSR}}
\newcommand{\GLS}{\text{GLS}}
\newcommand{\df}{\text{df}}
\newcommand{\ti}{\text{time}}
\newcommand{\outcome}{\text{outcome}}
\newcommand{\exposure}{\text{exposure}}
\newcommand{\NA}{\text{NA}}
\newcommand{\cp}{\text{cp}}
\newcommand{\tip}{\textbf{\faLightbulb}: }
\newcommand{\warning}{\textbf{\Biohazard}: }
\newcommand{\blank}[1][4cm]{\underline{\hspace{#1}}}

% Define theorem-style environment (italic body, bold header)
\newtheorem{theorem}{Theorem}
\newtheorem{lemma}{Lemma}
\newtheorem{proposition}{Proposition}
\newtheorem{corollary}{Corollary}

% Define definition-style environment (upright body, bold header)
\theoremstyle{definition}
\newtheorem{definition}{Definition}
\newtheorem{example}{Example}
\newtheorem{problem}{\faPen \ Problem}
\newtheorem{concept}{\faPen \ Concept Checker}
\newtheorem{challenge}{\faFire \ Challenge}
\newtheorem{activity}{\faDiceThree \ Activity}

% Define idea-style environment (unnumbered definition)
\newtheorem*{idea}{Idea}

% Define remark-style environment (upright body, italic header)
\theoremstyle{remark}
\newtheorem{remark}{Remark}

% Define environment for solution box
\newsavebox{\solutioncontentbox}
\NewEnviron{solutionbox}[1][Solution]{
  \sbox{\solutioncontentbox}{
    \begin{minipage}{\dimexpr\linewidth-2\fboxsep\relax}
      \BODY
    \end{minipage}
  }
  \ifsolutions
    \begin{tcolorbox}[
      title=#1,
      colback=gray!5,
      colframe=black,
      breakable
    ]
      \BODY
    \end{tcolorbox}
  \else
    \begin{tcolorbox}[
      title=#1,
      colback=white,
      colframe=black,
      breakable,
      height=\dimexpr\ht\solutioncontentbox+\dp\solutioncontentbox+1cm\relax
    ]
    \end{tcolorbox}
  \fi
}

% Define environment for code chunk
\definecolor{codegray}{RGB}{245,245,245}
\definecolor{commentgreen}{RGB}{0,128,0}
\definecolor{keywordblue}{RGB}{0,0,180}
\lstset{
    language=R,
    basicstyle=\ttfamily\small,
    backgroundcolor=\color{codegray},
    commentstyle=\color{commentgreen},
    keywordstyle=\color{black},
    stringstyle=\color{red!70!black},
    showstringspaces=false,
    frame=single,
    rulecolor=\color{gray!30},
    breaklines=true,
    columns=fullflexible,
    keepspaces=true,
    xleftmargin=0.5em,
    framexleftmargin=0.5em,
    aboveskip=0.8em,
    belowskip=0.8em
}

% Define environment for code output
\lstnewenvironment{routput}
{
    \lstset{
        language={},
        basicstyle=\ttfamily\small,
        backgroundcolor=\color{codegray},
        frame=single,
        rulecolor=\color{gray!30},
        breaklines=true,
        columns=fullflexible,
        keepspaces=true,
        showstringspaces=false,
        upquote=true,
        xleftmargin=0.5em,
        framexleftmargin=0.5em,
        literate={}
    }
}
{}

% Begin document
\begin{document}
\pagestyle{fancy}
\fancyhead[L]{STAT 139: Introduction to Linear Models}
\fancyhead[R]{Fall 2026}
\begin{center}
\bf \Large
Section 4: Normality Assumption and $t$-Test \\[0.5em]
\normalsize
STAT 139 Teaching Team
\end{center}

% Section: Introduction
\section{Introduction}

\subsection{Logistics}

Welcome! The goal of our weekly section is to review content from lecture, with emphasis on big-picture intuition, mathematical derivation, and hands-on practice.

Specifically, we aim for section to be interactive, well-paced, inclusive, and fun! Please do not hesitate to reach out if you have any questions/feedback!

\subsection{Office Hours}

$\bullet$ See Canvas for the most accurate information!

\subsection{Attendance}

$\bullet$ Please notify the Teaching Fellow (TF) if your attendance hasn't been taken!

% Section: Adding the Normality Assumption
\section{Adding the Normality Assumption}

\begin{idea}
We began with only four assumptions, but the fifth assumption of \textit{Normality} allows us to extend our inference, especially when we have a limited sample size and cannot appeal to the Central Limit Theorem.
\end{idea}

\subsection{Point Estimation for Parameters}

\begin{definition}[ELIH]
We begin with the four assumptions of \textit{existence}, \textit{linearity}, \textit{independence}, and \textit{homoskedasticity}: \[\Var[\varepsilon_i] < \infty,\] \[\E[Y_i] = \beta_0 + \beta_1 x_i,\] \[\varepsilon_i \indep \varepsilon_j,\] and \[\Var[\varepsilon_i] = \sigma^2\] for all $i, j \in \{1, \ldots, n\}$ and $i \neq j$.
\end{definition}

\begin{definition}[OLS estimators]
$(\hat{\beta}_{0, \OLS}, \hat{\beta}_{1, \OLS}) = \arg \min_{(\beta_0, \beta_1) \in \R^2} \left( \sum_{i = 1}^{n} (Y_i - (\beta_0 + \beta_1 x_i))^2 \right)$, which can be derived under ELIH: \[\hat{\beta}_{1, \OLS} = \frac{\sum_{i = 1}^{n} (x_i - \bar{x})(Y_i - \bar{Y})}{\sum_{i = 1}^{n} (x_i - \bar{x})^2}\] and \[\hat{\beta}_{0, \OLS} = \bar{Y} - \hat{\beta}_{1, \OLS} \bar{x}.\]
\end{definition}

\begin{definition}[Gauss-Markov Theorem]
Under ELIH, the OLS estimators are \textit{optimal} within the class of linear unbiased estimators. 

$\bullet$ I.e., they are BLUE: best linear unbiased estimators.
\end{definition}

\begin{definition}[ELIHN]
When we add \textit{normality}, the five main assumptions can be summarized as \[\varepsilon_i \simiid \N(0, \sigma^2),\] with $\sigma^2 < \infty$ for all $i \in \{1, \ldots, n\}$. Equivalently, \[Y_i \simind \N(\beta_0 + \beta_1 x_i, \sigma^2).\]

$\bullet$ Normality is \textit{not} needed for OLS point estimation but gives exact finite-sample distributions for \textit{hypothesis tests} and \textit{confidence intervals}.\footnote{Unless otherwise noted, we will assume ELIHN.}
\end{definition}

\begin{definition}[OLS-MLE connection]
Under ELIHN, the OLS estimators for $\beta_0$ and $\beta_1$ are the maximum likelihood estimators (MLEs): \[\hat{\beta}_{1, \MLE} = \frac{\sum_{i = 1}^{n} (x_i - \bar{x})(Y_i - \bar{Y})}{\sum_{i = 1}^{n} (x_i - \bar{x})^2}\] and \[\hat{\beta}_{0, \MLE} = \bar{Y} - \hat{\beta}_{1, \MLE} \bar{x}.\footnote{Unless otherwise noted, $\hat{\beta}_{1}$ refers to $\hat{\beta}_{1, \OLS} = \hat{\beta}_{1, \MLE}$ and $\hat{\beta}_{0}$ refers to $\hat{\beta}_{0, \OLS} = \hat{\beta}_{0, \MLE}$.}\]

$\bullet$ However, \[\hat{\sigma}^2_{\MLE} = \frac{1}{n} \sum_{i = 1}^{n} (Y_i - \hat{\beta}_0 - \hat{\beta}_1 x_i)^2,\] which is \textit{not} equal to \[\hat{\sigma}^2_{\OLS} = \frac{1}{n - 2} \sum_{i = 1}^{n} (Y_i - \hat{\beta}_0 - \hat{\beta}_1 x_i)^2.\footnote{See Appendix~\ref{app:ols-mle}.}\]
\end{definition}

\begin{concept}
Which estimator is unbiased: $\hat{\sigma}^2_{\MLE}$ or $\hat{\sigma}^2_{\OLS}$?
\end{concept}

\begin{solutionbox}
$\hat{\sigma}^2_{\OLS}$ is an unbiased estimator while $\hat{\sigma}^2_{\MLE}$ is not. Recall that $\frac{n - 2}{\sigma^2} \hat{\sigma}^2_{\OLS} \sim \chi^{2}_{n - 2}$. Thus, $\E[\frac{n - 2}{\sigma^2} \hat{\sigma}^2_{\OLS}] = n - 2$ by expectation of chi-squared, and $\E[\frac{n - 2}{\sigma^2} \hat{\sigma}^2_{\OLS}] = \frac{n - 2}{\sigma^2} \E[\hat{\sigma}^2_{\OLS}]$ by linearity, so $\E[\hat{\sigma}^2_{\OLS}] = \sigma^2$ by algebra.
\end{solutionbox}

\subsection{$t$-Based Inference for Parameters}

\begin{definition}[Distributions]
Under ELIHN, we have the following distributions for our estimators: \[\hat{\beta}_{1} \sim \N \left( \beta_1, \frac{\sigma^2}{s_{xx}} \right),\] \[\hat{\beta}_{0} \sim \N \left( \beta_0, \sigma^2 \left( \frac{1}{n} + \frac{\bar{x}^2}{s_{xx}} \right) \right),\] and \[\frac{n - 2}{\sigma^2} \hat{\sigma}^2_{\OLS} \sim \chi^2_{n - 2}.\]
\end{definition}

\begin{concept}
Assume that $\sigma^2$ is unknown. Can we use the distribution of $\hat{\beta}_{1}$ to find a $95\%$ confidence interval for $\beta_{1}$?
\end{concept}

\begin{solutionbox}
We know that $\hat{\beta}_{1} \sim \N \left( \beta_1, \frac{\sigma^2}{s_{xx}} \right)$, so if $\sigma^2$ were known, then we could construct a CI centered at $\hat{\beta}_1$. Using $z_{0.975} = 1.96$, we'd have $\hat{\beta}_{1} \pm 1.96 \times \sqrt{\frac{\sigma^2}{s_{xx}}}$. However, we don't know $\sigma^2$, which is why we must plug in an estimator like $\hat{\sigma}^2_{\OLS}$. This changes the distribution to $t$ instead of Normal!
\end{solutionbox}

\begin{definition}[Hypothesis test for slope]
To test $H_0: \beta_1 = 0$ vs. $H_A: \beta_1 \neq 0$, we use \[T = \frac{\hat{\beta}_{1} - 0}{\hat{\sigma}_{\OLS} \sqrt{\frac{1}{s_{xx}}}} \simnull t_{n - 2}\] as our \textit{test statistic} (i.e., its \textit{exact} distribution is known under $H_0$).
\end{definition}

\begin{concept}
Suppose that instead we wanted to test $H_0: \beta_1 = 139$ vs. $H_A: \beta_1 \neq 139$. What test statistic would we use? What would its distribution be under $H_0$?
\end{concept}

\begin{solutionbox}
To test $H_0: \beta_1 = 139$ vs. $H_A: \beta_1 \neq 139$, we would use $T = \frac{\hat{\beta}_{1} - 139}{\hat{\sigma}_{\OLS} \sqrt{\frac{1}{s_{xx}}}} \simnull t_{n - 2}$. Recall that $\hat{\beta}_{1} \sim \N \left( \beta_1, \frac{\sigma^2}{s_{xx}} \right)$, so $\frac{\hat{\beta}_{1} - \beta_1}{\sigma / \sqrt{s_{xx}}} \sim \N(0, 1)$ by standardizing. Additionally, $\frac{n - 2}{\sigma^2} \hat{\sigma}^2_{\OLS} \sim \chi^2_{n - 2}$, with $\frac{n - 2}{\sigma^2} \hat{\sigma}^2_{\OLS} \indep \frac{\hat{\beta}_{1} - \beta_1}{\sigma / \sqrt{s_{xx}}}$. Thus, for all values of $\beta_1$, $\frac{\hat{\beta}_{1} - \beta_1}{\sigma / \sqrt{s_{xx}}} \div \sqrt{\frac{n - 2}{\sigma^2} \hat{\sigma}^2_{\OLS} / (n - 2)} = \frac{\hat{\beta}_{1} - \beta_1}{\sigma / \sqrt{s_{xx}}} \div \frac{\hat{\sigma}_{\OLS}}{\sigma} = \frac{\hat{\beta}_{1} - \beta_1}{\hat{\sigma}_{\OLS} \sqrt{1 / s_{xx}}} \sim t_{n - 2}$ by the story of the $t$.
\end{solutionbox}


\begin{definition}[Confidence interval for slope]
Our $100 (1 - \alpha) \%$ confidence interval for $\beta_1$ is \[\hat{\beta}_{1} \pm t_{1 - \frac{\alpha}{2}, n - 2} \times \hat{\sigma}_{\OLS} \sqrt{1 / s_{xx}}.\]
\end{definition}

\begin{definition}[Hypothesis test for intercept]
To test $H_0: \beta_0 = 0$ vs. $H_A: \beta_0 \neq 0$, we use \[T = \frac{\hat{\beta}_{0} - 0}{\hat{\sigma}_{\OLS} \sqrt{\frac{1}{n} + \frac{\bar{x}^2}{s_{xx}}}} \simnull t_{n - 2}\] as our \textit{test statistic} (i.e., its \textit{exact} distribution is known under $H_0$).
\end{definition}

\begin{definition}[Confidence interval for intercept]
Our $100 (1 - \alpha) \%$ confidence interval for $\beta_0$ is \[\hat{\beta}_{0} \pm t_{1 - \frac{\alpha}{2}, n - 2} \times \hat{\sigma}_{\OLS} \sqrt{\frac{1}{n} + \frac{\bar{x}^2}{s_{xx}}}.\]
\end{definition}

\begin{concept}
We use $t_{n - 2}$ (and multiply by $\frac{1}{n - 2}$ in the formula for $\hat{\sigma}^{2}_{\OLS}$) because we lose $2$ degrees of freedom. Why? Suppose that we observe $y_1 = 1$ and $y_2 = 1$, with $n = 3$. If we tell you that $\bar{y} = 1$, then what must $y_3$ be?
\end{concept}

\begin{solutionbox}
Deterministically, $y_3 = 1$—it's no longer random, so we lose a degree of freedom. To calculate $\hat{\sigma}^{2}_{\OLS}$, notice that we must insert values for $\hat{\beta}_0$ and $\hat{\beta}_1$, so we're \enquote{double dipping} in the data, losing $2$ degrees of freedom.
\end{solutionbox}

\subsection{Point Estimation for Responses}

\begin{definition}[Point estimator]
Suppose that we're interested in the response $Y$ at some predictor value $X = x_0$. Our \textit{best guess} is \[\hat{\mu}_{Y \mid X = x_0} = \hat{\beta}_0 + \hat{\beta}_1 x_0.\]
\end{definition}

\subsection{$t$-Based Inference for Responses}

\begin{definition}[Population mean confidence interval]
We want to find a confidence interval for the \textit{population mean response} at some predictor value—i.e., $\mu_{Y \mid X = x_0} = \E[Y \mid X = x_0]$.

$\bullet$ E.g., what is the average height of children whose midparent height is 69.21 inches?

$\bullet$ Our $100 (1 - \alpha) \%$ confidence interval for $\mu_{Y \mid X = x_0}$ is \[(\hat{\beta}_{0} + \hat{\beta}_{1} x_0) \pm t_{1 - \frac{\alpha}{2}, n - 2} \times \hat{\sigma}_{\OLS} \sqrt{\frac{1}{n} + \frac{(x_0 - \bar{x})^2}{s_{xx}}}.\]
\end{definition}

\begin{definition}[Single observation prediction interval]
We want to find a prediction interval for a \textit{new response} at some predictor value—i.e., $Y_0 = Y \mid X = x_0$.

$\bullet$ E.g., what is the predicted height of a child whose midparent height is 69.21 inches?

$\bullet$ Our $100 (1 - \alpha) \%$ prediction interval for $Y_0$ is \[(\hat{\beta}_{0} + \hat{\beta}_{1} x_0) \pm t_{1 - \frac{\alpha}{2}, n - 2} \times \hat{\sigma}_{\OLS} \sqrt{1 + \frac{1}{n} + \frac{(x_0 - \bar{x})^2}{s_{xx}}}.\]

$\bullet$ We can decompose the prediction error: \[\underbrace{Y_0 - \hat{\mu}_{Y \mid X = x_0}}_{\text{prediction error}} = \underbrace{(Y_0 - \mu_{Y \mid X = x_0})}_{\text{random sampling error}} + \underbrace{(\mu_{Y \mid X = x_0} - \hat{\mu}_{Y \mid X = x_0})}_{\text{estimation error}}.\]
\end{definition}

\begin{concept}
Which has more uncertainty: a population mean confidence interval or a single observation prediction interval? 
\end{concept}

\begin{solutionbox}
The single observation prediction interval has more uncertainty, reflected by the additional $+ 1$ in the margin of error. This makes sense since there are so many other factors that go into a new response (e.g., for predicted height of a child, we could look at the nutrition/sleep/lifestyle of the child).
\end{solutionbox}

% Section: $t$-Test as a Linear Model
\section{$t$-Test as a Linear Model}

\begin{idea}
We can connect what we learned about the pooled two-sample $t$-test to the simple linear regression model!
\end{idea}

\begin{definition}[$t$-distribution]
Let $Z \sim \N(0, 1)$ and $Y \sim \chi^{2}_{\nu}$, with $Z \indep Y$. Then \[T = \frac{Z}{\sqrt{Y / \nu}} \sim t_{\nu}.\]
\end{definition}

\begin{definition}[One-sample $t$-test]
Let $Y_{1}, \ldots, Y_{n} \simiid \N(\mu, \sigma^2)$, with $\mu$ and $\sigma^2$ unknown. Define the \textit{sample variance} as \[S^2 = \frac{1}{n - 1} \sum_{i = 1}^{n} (Y_{i} - \bar{Y})^2.\] We want to conduct the two-sided test $H_{0}: \mu = \mu_{0}$ vs. $H_{A}: \mu \neq \mu_{0}$. We use \[T = \frac{\bar{Y} - \mu_{0}}{S / \sqrt{n}} \simnull t_{n - 1}\] as our \textit{test statistic} (i.e., its \textit{exact} distribution is known under $H_0$).
\end{definition}

\begin{definition}[Pooled two-sample $t$-test]
Let $Y_{1, 1}, \ldots, Y_{1, n_1} \simiid \N(\mu_{1}, \sigma^2)$ and $Y_{2, 1}, \ldots, Y_{2, n_2} \simiid \N(\mu_{2}, \sigma^2)$, with the two samples independent and $\mu_{1}$, $\mu_{2}$, and $\sigma^2$ unknown. Define the \textit{pooled sample variance} as \[S^2_p = \frac{(n_1 - 1) S^2_1 + (n_2 - 1) S^2_2}{n_1 + n_2 - 2}.\] We want to conduct the two-sided test $H_{0}: \mu_1 = \mu_2$ vs. $H_{A}: \mu_1 \neq \mu_2$. We use \[T = \frac{\bar{Y}_{1} - \bar{Y}_{2}}{S_p \sqrt{\frac{1}{n_1} + \frac{1}{n_2}}} \simnull t_{n_1 + n_2 - 2}\] as our \textit{test statistic} (i.e., its \textit{exact} distribution is known under $H_0$).
\end{definition}

\begin{definition}[Pooled two-sample $t$-test as a linear model]
We can express the setup as \[Y_{j, i} = \mu_{j} + \varepsilon_{j, i},\] where $\varepsilon_{j, i} \simiid \N(0, \sigma^2)$ for all $j \in \{1, 2\}$ and $i \in \{1, \ldots, n_j\}$.

$\bullet$ Equivalently, \[Y_{j, i} = \mu_{1} \indic\{\text{$j = 1$}\} + \mu_{2} \indic\{\text{$j = 2$}\} + \varepsilon_{j, i},\] which can be rewritten in matrix notation as

\[
\begin{pmatrix}
Y_{1, 1} \\
\vdots \\
Y_{1, n_1} \\
Y_{2, 1} \\
\vdots \\
Y_{2, n_2}
\end{pmatrix}
=
\begin{pmatrix}
1 & 0 \\
\vdots & \vdots \\
1 & 0 \\
0 & 1 \\
\vdots & \vdots \\
0 & 1
\end{pmatrix}
\begin{pmatrix}
\mu_{1} \\
\mu_{2}
\end{pmatrix}
+
\begin{pmatrix}
\varepsilon_{1, 1} \\
\vdots \\
\varepsilon_{1, n_1} \\
\varepsilon_{2, 1} \\
\vdots \\
\varepsilon_{2, n_2}
\end{pmatrix}.
\]

$\bullet$ As a linear model, let $\beta_0 = \mu_1$ and $\beta_1 = \mu_2 - \mu_1$ such that \[Y_{j, i} = \beta_{0} + \beta_{1} \indic\{\text{$j = 2$}\} + \varepsilon_{j, i},\] which can be rewritten in matrix notation as

\[
\begin{pmatrix}
Y_{1, 1} \\
\vdots \\
Y_{1, n_1} \\
Y_{2, 1} \\
\vdots \\
Y_{2, n_2}
\end{pmatrix}
=
\begin{pmatrix}
1 & 0 \\
\vdots & \vdots \\
1 & 0 \\
1 & 1 \\
\vdots & \vdots \\
1 & 1
\end{pmatrix}
\begin{pmatrix}
\beta_{0} \\
\beta_{1}
\end{pmatrix}
+
\begin{pmatrix}
\varepsilon_{1, 1} \\
\vdots \\
\varepsilon_{1, n_1} \\
\varepsilon_{2, 1} \\
\vdots \\
\varepsilon_{2, n_2}
\end{pmatrix}.
\]
\end{definition}

% Section: Exercises
\section{Exercises}

\begin{problem}
For this problem, refer to the code and output below.

(a) Walk through the code. Specifically, it may help to fill in the comments.

(b) What are the formulas for \texttt{Estimate:(Intercept) = 138.902} and \texttt{Estimate:indic = -29.236} from \texttt{lm()}?

(c) Why is \texttt{t value:indic = -118.8} from \texttt{lm()} different from \texttt{t = 115.17} from \texttt{t.test()}? 
\end{problem}

\begin{lstlisting}
```{r}
# 1: 
set.seed(139)

# 2: 
n_1 = 30
n_2 = 40
mu_1 = 139
mu_2 = 110
sigma = 1

# 3: 
Y_1_vec <- rnorm(n_1, mu_1, sigma)
Y_2_vec <- rnorm(n_2, mu_2, sigma)

# 4: 
Y <- c(Y_1_vec, Y_2_vec)
indic <- c(rep(0, n_1), rep(1, n_2))
df <- as.data.frame(cbind(Y, indic))

# 5: 
t.test(Y_1_vec, Y_2_vec)

# 6: 
summary(lm(Y ~ indic, data = df))
```
\end{lstlisting}

\begin{routput}
	Welch Two Sample t-test

data:  Y_1_vec and Y_2_vec
t = 115.17, df = 54.401, p-value < 2.2e-16
alternative hypothesis: true difference in means is not equal to 0
95 percent confidence interval:
 28.72724 29.74491
sample estimates:
mean of x mean of y 
  138.902   109.666 

Call:
lm(formula = Y ~ indic, data = df)

Coefficients:
            Estimate Std. Error t value Pr(>|t|)    
(Intercept)  138.902      0.186   746.8   <2e-16 ***
indic        -29.236      0.246  -118.8   <2e-16 ***
---
\end{routput}

\begin{solutionbox}
(a) 
\# 1: Set seed for reproducibility
\# 2: Define parameters
\# 3: Generate random data
\# 4: Define df for lm()
\# 5: Run t-test
\# 6: Run linear regression

(b) For \texttt{Estimate:(Intercept) = 138.902}, this is $\bar{y}_1$. For \texttt{Estimate:indic = -29.236}, this is $\bar{y}_2 - \bar{y}_1$. In general for linear regression with a categorical predictor in \texttt{lm()}, the estimate of the $y$-intercept is the sample mean for the baseline group (here, $\bar{y}_1$) and the estimate of the indicator for group $j$ is the sample mean for group $j$ minus the sample mean for the baseline group (here, $\bar{y}_2 - \bar{y}_1$).

(c) Importantly, \texttt{t.test()}, by default, uses an unpooled $t$-test while the standard \texttt{lm()} fit assumes a common variance across the two groups and therefore yields the equivalent of the pooled two-sample $t$-test (which is appropriate here since the two groups do have shared variance). If we instead ran \texttt{t.test(Y\_1\_vec, Y\_2\_vec, var.equal = TRUE)}, then we'd see \texttt{t = 118.83}. As explained in (b), \texttt{lm()} flips the order, so \texttt{t = 118.83} from \texttt{t.test()} corresponds to \texttt{t value:indic = -118.8} from \texttt{lm()}.
\end{solutionbox}

\begin{problem}
Let's derive the population mean confidence interval and single observation prediction interval, which are both centered around $\hat{\mu}_{Y \mid X = x_0} = \hat{\beta}_0 + \hat{\beta}_1 x_0$.

(a) Find $\E[\hat{\mu}_{Y \mid X = x_0}]$.

(b) Find $\Cov[\hat{\beta}_0, \hat{\beta}_1]$.

(c) Find $\Var[\hat{\mu}_{Y \mid X = x_0}]$.

(d) Find the distribution of $\hat{\mu}_{Y \mid X = x_0}$.

(e) Recall that $Y_i \simind \N(\beta_0 + \beta_1 x_i, \sigma^2)$. Thus, $Y_0 \sim \N(\beta_0 + \beta_1 x_0, \sigma^2)$. This is a new response, so $Y_0$ is independent of $Y_1, \ldots, Y_n$—i.e., the data used to calculate $\hat{\beta}_0$, $\hat{\beta}_1$, and $\hat{\sigma}^2_{\OLS}$. Recall that the prediction error is $Y_0 - \hat{\mu}_{Y \mid X = x_0}$, whose expectation is $0$. Find $\Var[Y_0 - \hat{\mu}_{Y \mid X = x_0}]$.
\end{problem}

\begin{solutionbox}
{
\scriptsize
(a)
\[
\begin{aligned}
\E[\hat{\mu}_{Y \mid X = x_0}] &= \E[\hat{\beta}_0 + \hat{\beta}_1 x_0] && \text{by substituting} \\
&= \E[\hat{\beta}_0] + x_0 \E[\hat{\beta}_1] && \text{by linearity} \\
&= \beta_0 + \beta_1 x_0 && \text{by unbiasedness of OLS estimators}
\end{aligned}
\]

(b)
\[
\begin{aligned}
\Cov[\hat{\beta}_0, \hat{\beta}_1] &= \Cov[\bar{Y} - \hat{\beta}_{1} \bar{x}, \hat{\beta}_{1}] && \text{by substituting} \\
&= \Cov[\bar{Y}, \hat{\beta}_{1}] + \Cov[- \hat{\beta}_{1} \bar{x}, \hat{\beta}_{1}] && \text{by distributing} \\
&= \Cov[\bar{Y}, \hat{\beta}_{1}] - \bar{x} \Var[\hat{\beta}_{1}] && \text{by bilinearity} \\
&= \Cov \left[ \frac{1}{n} \sum_{i = 1}^{n} Y_i, \frac{\sum_{j = 1}^{n} (x_j - \bar{x})(Y_j - \bar{Y})}{s_{xx}} \right] - \bar{x} \Var[\hat{\beta}_{1}] && \text{by substituting} \\
&= \Cov \left[ \frac{1}{n} \sum_{i = 1}^{n} Y_i, \frac{\sum_{j = 1}^{n} (x_j - \bar{x}) Y_j}{s_{xx}} \right] - \bar{x} \Var[\hat{\beta}_{1}] && \text{since $\sum_{i = 1}^{n} (x_i - \bar{x}) (Y_i - \bar{Y}) = \sum_{i = 1}^{n} (x_i - \bar{x}) Y_i$} \\
&= \left( \frac{1}{n s_{xx}} \right) \Cov \left[ \sum_{i = 1}^{n} Y_i, \sum_{j = 1}^{n} (x_j - \bar{x}) Y_j \right] - \bar{x} \Var[\hat{\beta}_{1}] && \text{by bilinearity} \\
&= \left( \frac{1}{n s_{xx}} \right) \sum_{i = 1}^{n} \sum_{j = 1}^{n} \Cov \left[ Y_i, (x_j - \bar{x}) Y_j \right] - \bar{x} \Var[\hat{\beta}_{1}] && \text{by distributing} \\
&= \left( \frac{1}{n s_{xx}} \right) \sum_{i = 1}^{n} \sum_{j = 1}^{n} (x_j - \bar{x}) \Cov \left[ Y_i, Y_j \right] - \bar{x} \Var[\hat{\beta}_{1}] && \text{by bilinearity} \\
&= \left( \frac{1}{n s_{xx}} \right) \left( \sum_{i = j} (x_j - \bar{x}) \Cov \left[ Y_i, Y_j \right] + \sum_{i \neq j} (x_j - \bar{x}) \Cov \left[ Y_i, Y_j \right] \right) - \bar{x} \Var[\hat{\beta}_{1}] && \text{by separating summation} \\
&= \left( \frac{1}{n s_{xx}} \right) \left( \sigma^2 \sum_{i = j} (x_j - \bar{x}) + 0 \sum_{i \neq j} (x_j - \bar{x}) \right) - \bar{x} \Var[\hat{\beta}_{1}] && \text{by independence} \\
&= \left( \frac{1}{n s_{xx}} \right) \left( \sigma^2 0 \right) - \bar{x} \Var[\hat{\beta}_{1}] && \text{since $\sum_{i = 1}^{n} (x_i - \bar{x}) = 0$} \\
&= - \bar{x} \Var[\hat{\beta}_{1}] && \text{by simplifying}
\end{aligned}
\]

(c)
\[
\begin{aligned}
\Var[\hat{\mu}_{Y \mid X = x_0}] &= \Var[\hat{\beta}_0 + \hat{\beta}_1 x_0] && \text{by substituting} \\
&= \Var[\hat{\beta}_0] + x_0^2 \Var[\hat{\beta}_1] + 2 x_0 \Cov[\hat{\beta}_0, \hat{\beta}_1] && \text{by bilinearity} \\
&= \Var[\hat{\beta}_0] + x_0^2 \Var[\hat{\beta}_1] + 2 x_0 (- \bar{x} \Var[\hat{\beta}_{1}]) && \text{by substituting} \\
&= \Var[\hat{\beta}_0] + \Var[\hat{\beta}_1] (x_0^2 - 2 x_0 \bar{x}) && \text{by simplifying} \\
&= \sigma^2 \left( \frac{1}{n} + \frac{\bar{x}^2}{s_{xx}} \right) + \frac{\sigma^2}{s_{xx}} (x_0^2 - 2 x_0 \bar{x}) && \text{by substituting} \\
&= \sigma^2 \left( \frac{1}{n} + \frac{(x_0 - \bar{x})^2}{s_{xx}} \right) && \text{by simplifying}
\end{aligned}
\]

(d) $\hat{\mu}_{Y \mid X = x_0} = \hat{\beta}_0 + \hat{\beta}_1 x_0$, where $(\hat{\beta}_0, \hat{\beta}_1)$ is Multivariate Normal, so $\hat{\mu}_{Y \mid X = x_0}$ is Normal as a linear combination. Thus, $\hat{\mu}_{Y \mid X = x_0} \sim \N \left( \beta_0 + \beta_1 x_0, \sigma^2 \left( \frac{1}{n} + \frac{(x_0 - \bar{x})^2}{s_{xx}} \right) \right)$.

(e)
\[
\begin{aligned}
\Var[Y_0 - \hat{\mu}_{Y \mid X = x_0}] &= \Var[Y_0] + \Var[\hat{\mu}_{Y \mid X = x_0}] - 2 \Cov[Y_0, \hat{\mu}_{Y \mid X = x_0}] && \text{by bilinearity} \\
&= \Var[Y_0] + \Var[\hat{\mu}_{Y \mid X = x_0}] && \text{by independence} \\
&= \sigma^2 + \sigma^2 \left( \frac{1}{n} + \frac{(x_0 - \bar{x})^2}{s_{xx}} \right) && \text{by substituting} \\
&= \sigma^2 \left(1 + \frac{1}{n} + \frac{(x_0 - \bar{x})^2}{s_{xx}} \right) && \text{by substituting}
\end{aligned}
\]
}
\end{solutionbox}

% Appendix
\appendix
\section{Proofs}

\subsection{OLS-MLE Connection}
\label{app:ols-mle}

\begin{proof}
Assume ELIHN. We want to show that $\hat{\beta}_{1, \MLE} = \hat{\beta}_{1, \OLS}$, $\hat{\beta}_{0, \MLE} = \hat{\beta}_{0, \OLS}$, and $\hat{\sigma}^2_{\MLE} \neq \hat{\sigma}^2_{\OLS}$. Recall that $\hat{\beta}_{1, \OLS} = \frac{\sum_{i = 1}^{n} (x_i - \bar{x})(Y_i - \bar{Y})}{\sum_{i = 1}^{n} (x_i - \bar{x})^2}$, $\hat{\beta}_{0, \OLS} = \bar{Y} - \hat{\beta}_{1, \OLS} \bar{x}$, and $\hat{\sigma}^2_{\OLS} = \frac{1}{n - 2} \sum_{i = 1}^{n} (Y_i - \hat{\beta}_0 - \hat{\beta}_1 x_i)^2$.

By definition, the maximum likelihood estimators maximize the likelihood function—i.e., $\widehat{\vec{\theta}}_{\MLE} = \arg \max_{\vec{\theta} \in \Theta} \Like(\vec{\theta}; \vec{Y})$. Here, $\vec{\theta} = (\beta_0, \beta_1, \sigma^2)$, and $Y_i \simind \N(\beta_0 + \beta_1 x_i, \sigma^2)$.

\[
\begin{aligned}
\Like(\beta_0, \beta_1, \sigma^2; \vec{Y}) &= f_{\vec{Y}}(\vec{Y}; \beta_0, \beta_1, \sigma^2) && \text{by def. of likelihood} \\
&= \prod_{i = 1}^{n} f_{Y_i}(Y_i; \beta_0, \beta_1, \sigma^2) && \text{by independence} \\
&= \prod_{i = 1}^{n} \frac{1}{\sqrt{2 \pi \sigma^2}} \exp \left(
- \frac{1}{2 \sigma^2} (Y_i - \beta_0 - \beta_1 x_i)^2 \right) && \text{by Normal PDF} \\
&= \frac{1}{(2 \pi \sigma^2)^{\frac{n}{2}}} \exp \left( - \frac{1}{2 \sigma^2} \sum_{i = 1}^{n} (Y_i - \beta_0 - \beta_1 x_i)^2 \right) && \text{by multiplying}
\end{aligned}
\]

\[
\begin{aligned}
\ell(\beta_0, \beta_1, \sigma^2; \vec{Y}) &= \log \Like(\beta_0, \beta_1, \sigma^2; \vec{Y}) && \text{by def. of log-likelihood} \\
&= - \frac{n}{2}\log(2\pi) - \frac{n}{2}\log(\sigma^2) - \frac{1}{2\sigma^2} \sum_{i = 1}^{n} (Y_i - \beta_0 - \beta_1 x_i)^2 && \text{by simplifying}
\end{aligned}
\]

For any fixed $\sigma^2 > 0$, the first two terms do not depend on $\beta_0$ or $\beta_1$, and $\frac{1}{2\sigma^2} > 0$. Thus, \textit{maximizing} the log-likelihood with respect to $(\beta_0,\beta_1)$ is equivalent to \textit{minimizing} \[\sum_{i = 1}^{n} (Y_i - \beta_0 - \beta_1 x_i)^2,\] which is exactly the OLS objective. Therefore, \[\hat{\beta}_{1,\MLE} = \hat{\beta}_{1,\OLS}\] and \[\hat{\beta}_{0,\MLE} = \hat{\beta}_{0,\OLS}.\]

Finally, let's maximize $\ell(\beta_0, \beta_1, \sigma^2; \vec{Y})$ with respect to $\sigma^2$. Let $c = \sum_{i = 1}^{n} (Y_i - \beta_0 - \beta_1 x_i)^2$, which does not depend on $\sigma^2$.

\[
\begin{aligned}
\frac{\partial}{\partial \sigma^2} \ell(\beta_0, \beta_1, \sigma^2; \vec{Y}) &= \frac{\partial}{\partial \sigma^2} \left( - \frac{n}{2}\log(2\pi) - \frac{n}{2}\log(\sigma^2) - \frac{1}{2\sigma^2} c \right) && \text{by substituting} \\
&= - \frac{n}{2} \times \frac{\partial}{\partial \sigma^2} \left( \log(\sigma^2) \right) - \frac{c}{2} \times \frac{\partial}{\partial \sigma^2} \left( \frac{1}{\sigma^2} \right) && \text{by linearity} \\
&= - \frac{n}{2 \sigma^2} + \frac{c}{2 (\sigma^2)^2} && \text{by simplifying}
\end{aligned}
\]

\[
\begin{aligned}
0 = - \frac{n}{2 \sigma^2} + \frac{c}{2 (\sigma^2)^2} &\implies \frac{n}{2 \sigma^2} = \frac{c}{2 (\sigma^2)^2} && \text{by adding} \\
&\implies n 2 (\sigma^2)^2 = 2 \sigma^2 c && \text{by cross-multiplying} \\
&\implies \sigma^2 = \frac{1}{n} \sum_{i = 1}^{n} (Y_i - \beta_0 - \beta_1 x_i)^2 && \text{by simplifying} \\
\end{aligned}
\]

Therefore, \[\hat{\sigma}^2_{\MLE} = \frac{1}{n} \sum_{i = 1}^{n} (Y_i - \hat{\beta}_0 - \hat{\beta}_1 x_i)^2,\] which is \textit{not} equal to \[\hat{\sigma}^2_{\OLS} = \frac{1}{n - 2} \sum_{i = 1}^{n} (Y_i - \hat{\beta}_0 - \hat{\beta}_1 x_i)^2.\]
\end{proof}

\end{document}